\chapter{Exact inference for the Normalised Power Prior}
\label{sec:exact}

In this chapter, we develop exact Markov chain Monte Carlo methods for inference under the normalised power prior.

The full joint posterior of interest is
\begin{equation}
    \pi(\theta, \eta \mid D, D_0)
    \propto
    L(D \mid \theta)\,
    \frac{L_0(D_0 \mid \theta)^{\eta}\, \pi_0(\theta)}{c_0(\eta)}\,
    \pi_A(\eta),
    \label{eq:full-posterior-exact}
\end{equation}
and the central challenge is the intractable ratio $c_0(\eta)/c_0(\eta')$ that appears in any Metropolis--Hastings update of $\eta$.

We address this challenge through three complementary constructions: the exchange algorithm (Section~\ref{sec:exact-exchange}), the Poisson estimator (Section~\ref{subsec:poisson-estimator}), and a pseudo-marginal approach based on importance sampling (Section~\ref{sec:pseudo-marginal-is}).

Updating $\theta$ given $(\eta, D, D_0)$ is straightforward: the full conditional
\begin{equation}
    \pi(\theta \mid \eta, D, D_0) \propto L(D \mid \theta)\, L_0(D_0 \mid \theta)^{\eta}\, \pi_0(\theta)
    \label{eq:theta-full-conditional}
\end{equation}
does not involve $c_0(\eta)$, so standard Metropolis--Hastings (Section~\ref{subsec:metropolis}) applies directly.

For a symmetric proposal $q(\theta' \mid \theta) = q(\theta \mid \theta')$, the acceptance ratio reduces to
\begin{equation}
    \alpha(\theta, \theta' \mid \eta)
    =
    \min\!\left\{1,\,
    \frac{L(D \mid \theta')\,\pi_0(\theta')}{L(D \mid \theta)\,\pi_0(\theta)}
    \left(\frac{L_0(D_0 \mid \theta')}{L_0(D_0 \mid \theta)}\right)^{\!\eta}
    \right\}.
\end{equation}

The remainder of this chapter focuses on the harder problem of updating $\eta$.

\section{Exchange Metropolis--Hastings}
\label{sec:exact-exchange}

Applying the exchange algorithm of Section~\ref{subsec:exchange-algorithm} to the NPP, the intractable quantity is $c_0(\eta)$ as defined in \eqref{eq:normalising-constant}.
\begin{equation}
c_0(\eta)
=
\int_\Theta L_0(D_0 \mid \theta)^\eta \pi_0(\theta) \dd \theta.
\end{equation}

If the current value of $\theta$ is fixed and a proposal $\eta'$ is generated from $q(\eta' \mid \eta)$, the exchange update draws
\begin{equation}
\theta^\star
\sim
\pi_{\eta'}(\theta \mid D_0)
=
\frac{L_0(D_0 \mid \theta)^{\eta'}\pi_0(\theta)}
{c_0(\eta')}.
\end{equation}
The exchange ratio for the $\eta$ update is then
\begin{align}
R_{\mathrm{ex}}(\eta,\eta';\theta,\theta^\star)
&=
\frac{\pi_A(\eta')}{\pi_A(\eta)}
\frac{q(\eta \mid \eta')}{q(\eta' \mid \eta)}
\frac{L_0(D_0 \mid \theta)^{\eta'-\eta}}
{L_0(D_0 \mid \theta^\star)^{\eta'-\eta}}.
\end{align}
The ratio $c_0(\eta)/c_0(\eta')$ has disappeared.
This cancellation is the main reason why the exchange algorithm is attractive for the normalised power prior.

According to the exchange algorithm framework of \citeonline{murray2012}, the auxiliary draw from $\pi_{\eta'}(\theta \mid D_0)$ introduces a factor $1/c_0(\eta')$ into the augmented target that exactly cancels the intractable ratio (see Appendix~\ref{proof:exchange-correctness} for a self-contained argument).

The exactness of the exchange algorithm depends on the ability to sample $\theta^\star$ exactly from $\pi_{\eta'}(\theta \mid D_0)$.
This assumption is strong.
It is satisfied in conjugate Gaussian models and in some conditionally conjugate models.
It is not automatically satisfied in general regression models or in non-conjugate likelihoods.
If $\theta^\star$ is generated by an inner MCMC chain that is not run to stationarity, the resulting algorithm is no longer exactly the exchange algorithm.
In that case, the method becomes an approximate exchange or double Metropolis--Hastings algorithm.

\section{Barker's method via Poisson estimator}

Updating $\eta$ is more delicate because the full conditional
\begin{equation}
    \pi(\eta \mid \theta, D, D_0) \propto \frac{\pi_A(\eta)}{c_0(\eta)}\, L_0(D_0 \mid \theta)^{\eta}
    \label{eq:eta-full-conditional}
\end{equation}
involves $c_0(\eta)$ explicitly.

The Metropolis--Hastings ratio for a proposal $\eta \to \eta'$ is
\begin{equation}
    \frac{\pi(\eta' \mid \theta)}{\pi(\eta \mid \theta)}
    =
    \frac{\pi_A(\eta')}{\pi_A(\eta)}
    \cdot
    \frac{c_0(\eta)}{c_0(\eta')}
    \cdot
    L_0(D_0 \mid \theta)^{\eta' - \eta},
\end{equation}
and the intractable ratio $c_0(\eta)/c_0(\eta')$ prevents direct evaluation.

Recalling the thermodynamic identity established in \eqref{eq:normalising-constant-log}, integrating from $\eta$ to $\eta'$ gives
\begin{equation}
    \label{eq:log-ratio-integral}
    \Delta(\eta, \eta')
    \coloneqq
    \log\frac{c_0(\eta')}{c_0(\eta)}
    =
    \int_{\eta}^{\eta'}
    \bE_{\pi_\tau}\bigl[\log L_0(D_0 \mid \theta)\bigr]
    \,\dd\tau,
\end{equation}
where $\pi_\tau(\theta) \propto L_0(D_0 \mid \theta)^\tau \pi_0(\theta)$ is the power prior at temperature $\tau$.

The goal of the following sections is to construct a non-negative unbiased estimator of $\exp\{-\Delta(\eta,\eta')\} = c_0(\eta)/c_0(\eta')$.


\subsection{Producing an unbiased estimator via a Poisson estimator}
\label{subsec:poisson-estimator}

Let $\pi_\tau$ denote the power posterior associated with the historical data.
For $\tau\in[0,1]$, define
\begin{equation}
    \pi_\tau(\theta)
    =
    \frac{L_0(D_0\mid \theta)^\tau \pi_0(\theta)}
    {c_0(\tau)},
    \qquad
    c_0(\tau)
    =
    \int_\Theta L_0(D_0\mid t)^\tau \,\dd P_0(t).
\end{equation}
For two values $\eta,\eta'\in[0,1]$, define
\begin{equation}
    \Delta(\eta,\eta')
    =
    \log\frac{c_0(\eta')}{c_0(\eta)}.
\end{equation}
Under the usual regularity conditions allowing differentiation under the integral sign, the thermodynamic identity gives
% Identity is proved in Appendix~\ref{proof:deriv-normalising-constant}.
\begin{equation}
    \Delta(\eta,\eta')
    =
    \int_{\eta}^{\eta'}
    \bE_{\theta\sim \pi_\tau}
    \left[
        \log L_0(D_0\mid \theta)
    \right]
    \,\dd \tau.
\end{equation}
Equivalently, if $U\sim\mathrm{Unif}(0,1)$ and $\tau_{\eta,\eta'}(U)=\eta+U(\eta'-\eta)$, then
\begin{equation}
    \Delta(\eta,\eta')
    =
    \bE_{U}
    \left[
        (\eta'-\eta)
        \bE_{\theta\sim \pi_{\tau_{\eta,\eta'}(U)}}
        \{\log L_0(D_0\mid \theta)\}
    \right].
\end{equation}

This representation suggests a proposal-specific random variable.
Let $\mathsf P_{\eta,\eta'}$ be the probability measure on $[0,1]\times\Theta$ obtained by first drawing $U\sim\mathrm{Unif}(0,1)$ and then drawing $\theta\sim\pi_{\tau_{\eta,\eta'}(U)}$.
Equivalently,
\begin{equation}
    \mathsf P_{\eta,\eta'}(\dd u,\dd\theta)
    =
    \dd u\,
    \pi_{\tau_{\eta,\eta'}(u)}(\dd\theta).
\end{equation}
Define
\begin{equation}
    Z_{\eta,\eta'}(U,\theta)
    =
    (\eta'-\eta)\log L_0(D_0\mid \theta).
\end{equation}
Then
\begin{equation}
    \bE_{\mathsf P_{\eta,\eta'}}[Z_{\eta,\eta'}]
    =
    \Delta(\eta,\eta').
\end{equation}

The objective is to construct an unbiased estimator of the ratio $c_0(\eta)/c_0(\eta')$.
This ratio can be written as
\begin{equation}
    \exp\{-\Delta(\eta,\eta')\}
    =
    \frac{c_0(\eta)}{c_0(\eta')}.
\end{equation}
Assume that there exist deterministic constants $a_{\eta,\eta'}$ and $b_{\eta,\eta'}$ such that
\begin{equation}
    a_{\eta,\eta'}
    \le
    Z_{\eta,\eta'}(U,\theta)
    \le
    b_{\eta,\eta'}
    \qquad
    \mathsf P_{\eta,\eta'}\text{-almost surely}.
\end{equation}
Set
\begin{equation}
    c_{\eta,\eta'}
    =
    b_{\eta,\eta'},
    \qquad
    \lambda_{\eta,\eta'}
    =
    b_{\eta,\eta'}-a_{\eta,\eta'}.
\end{equation}
Draw $K\sim\mathrm{Poisson}(\lambda_{\eta,\eta'})$.
Conditionally on $K$, draw
\begin{equation}
    (U_j,\theta_j)_{j=1}^K
    \overset{\mathrm{iid}}{\sim}
    \mathsf P_{\eta,\eta'}.
\end{equation}
The Poisson estimator is
\begin{equation}
    \widehat I(\eta,\eta')
    =
    \exp\{\lambda_{\eta,\eta'}-c_{\eta,\eta'}\}
    \prod_{j=1}^{K}
    \frac{
        c_{\eta,\eta'} -
        Z_{\eta,\eta'}(U_j,\theta_j)
    }{
        \lambda_{\eta,\eta'}
    }.
    \label{eq:poisson-estimator}
\end{equation}
By construction, each factor satisfies
\begin{equation}
    0
    \le
    \frac{
        c_{\eta,\eta'} -
        Z_{\eta,\eta'}(U_j,\theta_j)
    }{
        \lambda_{\eta,\eta'}
    }
    \le
    1.
\end{equation}
Therefore $\widehat I(\eta,\eta')$ is non-negative.
\begin{proof}
    
Using the exponential formula for a Poisson random variable, we obtain
\begin{align}
    \bE[\widehat I(\eta,\eta')]
    &=
    \exp\{\lambda_{\eta,\eta'}-c_{\eta,\eta'}\}
    \bE\left[
        \left\{
        \bE_{\mathsf P_{\eta,\eta'}}
        \left[
        \frac{c_{\eta,\eta'}-Z_{\eta,\eta'}}{\lambda_{\eta,\eta'}}
        \right]
        \right\}^K
    \right] \\
    &=
    \exp\{\lambda_{\eta,\eta'}-c_{\eta,\eta'}\}
    \exp\left\{
        \lambda_{\eta,\eta'}
        \left(
        \frac{c_{\eta,\eta'}-\Delta(\eta,\eta')}
             {\lambda_{\eta,\eta'}}
        -1
        \right)
    \right\} \\
    &=
    \exp\{-\Delta(\eta,\eta')\}.
\end{align}
Thus, $\widehat I(\eta,\eta')$ is an unbiased estimator of the normalising-constant ratio $c_0(\eta)/c_0(\eta')$.
\end{proof}

The efficiency of the estimator depends strongly on the envelope width $\lambda_{\eta,\eta'}=b_{\eta,\eta'}-a_{\eta,\eta'}$.
Large envelopes lead to large Poisson counts and high estimator variability.
Therefore, the practical usefulness of the construction depends on obtaining deterministic but sufficiently tight proposal-specific bounds for $Z_{\eta,\eta'}$.

It is important to distinguish this unbiased estimator from an exact Barker accept--reject mechanism.
Although $\widehat I(\eta,\eta')$ is unbiased for $c_0(\eta)/c_0(\eta')$, substituting a noisy estimate of the likelihood ratio directly into the Barker probability does not generally produce an exact Barker transition.
Exact Barker implementations require simulating the accept--reject event itself with the correct marginal probability, typically through a Bernoulli-factory or two-coin construction.

\subsubsection{Practical implementation}

The integral $\Delta(\eta,\eta')$ in \eqref{eq:log-ratio-integral} is estimated by Monte Carlo.

Drawing $\tau_k \overset{\mathrm{iid}}{\sim} \mathrm{Unif}(\eta,\eta')$ and, conditionally, $\widetilde{\theta}_{m,k} \sim \pi_{\tau_k}$ for $m = 1,\ldots,M$ and $k = 1,\ldots,K$, gives the double Monte Carlo estimator
\begin{equation}
    \widehat{\Delta}(\eta, \eta')
    =
    \frac{(\eta' - \eta)}{MK}
    \sum_{k=1}^{K} \sum_{m=1}^{M}
    \log L_0\bigl(D_0 \mid \widetilde{\theta}_{m,k}\bigr).
\end{equation}
This is an unbiased estimator of $\Delta(\eta,\eta')$ by the law of total expectation.

In the single-draw case ($K = M = 1$), this reduces to
\begin{equation}
    \label{eq:single-draw-delta}
    \widehat{\Delta}(\eta, \eta') = (\eta' - \eta)\,\log L_0(D_0 \mid \widetilde{\theta}_\star),
    \qquad
    \tau_\star \sim \mathrm{Unif}(\eta,\eta'),
    \quad
    \widetilde{\theta}_\star \sim \pi_{\tau_\star},
\end{equation}
which can be used in the Poisson estimator construction of Section~\ref{subsec:poisson-estimator}.

Note that only the mean of $Z_{\eta,\eta'}$, not its variance, enters the unbiasedness argument; the variance governs the efficiency of the estimator through the expected Poisson count $\lambda_{\eta,\eta'}$.

\subsection{Local envelopes for the Poisson estimator}
\label{subsec:local-envelopes}

The preceding construction does not require a single global envelope that is valid uniformly over all possible proposals.
Instead, for each proposed move $\eta\to\eta'$, it is enough to obtain deterministic bounds for the proposal-specific random variable
\begin{equation}
    Z_{\eta,\eta'}(U,\theta)
    =
    (\eta'-\eta)\log L_0(D_0\mid \theta),
    \qquad
    (U,\theta)\sim\mathsf P_{\eta,\eta'}.
\end{equation}
Suppose that the historical log-likelihood satisfies
\begin{equation}
    \ell_{\min}
    \le
    \log L_0(D_0\mid\theta)
    \le
    \ell_{\max}
    \qquad
    \text{for all }\theta\in\Theta.
\end{equation}
Then, for a fixed proposal $\eta\to\eta'$, define
\begin{equation}
    a_{\eta,\eta'}
    =
    \min\{(\eta'-\eta)\ell_{\min},
          (\eta'-\eta)\ell_{\max}\}.
\end{equation}
Similarly, define
\begin{equation}
    b_{\eta,\eta'}
    =
    \max\{(\eta'-\eta)\ell_{\min},
          (\eta'-\eta)\ell_{\max}\}.
\end{equation}
The resulting envelope width is
\begin{equation}
    b_{\eta,\eta'}-a_{\eta,\eta'}
    =
    |\eta'-\eta|(\ell_{\max}-\ell_{\min}).
\end{equation}
Hence a natural local choice is
\begin{equation}
    c_{\eta,\eta'}
    =
    b_{\eta,\eta'},
    \qquad
    \lambda_{\eta,\eta'}
    =
    |\eta'-\eta|(\ell_{\max}-\ell_{\min}).
\end{equation}
This choice makes the expected number of Poisson factors proportional to the proposal distance $|\eta'-\eta|$.
Consequently, local random-walk proposals in $\eta$ lead to smaller expected computational cost than large independent proposals.

For unbounded parameter spaces, such as Gaussian models with $\Theta=\mathds{R}^d$, the lower bound $\ell_{\min}$ is typically $-\infty$.
In that case, the above exact Poisson construction is not directly applicable without additional structure.
One possible solution is to work with a compactly supported prior.
For example, define
\begin{equation}
    \pi_{0,B}(\theta)
    \propto
    \pi_0(\theta)\mathbf 1_{\Theta_B}(\theta),
\end{equation}
where $\Theta_B\subset\Theta$ is compact.
Then deterministic bounds for $\log L_0(D_0\mid\theta)$ can be derived on $\Theta_B$.
In the Gaussian historical likelihood with known covariance, the log-likelihood is a concave quadratic function of $\theta$.
Its maximum is attained at the historical maximum likelihood estimator.
Its minimum over a compact box is attained on the boundary.
These deterministic bounds yield a valid local envelope for the Poisson estimator.

\subsection{Divide-and-conquer factorization}
\label{subsec:divide-and-conquer}


The computational difficulty of the Poisson estimator increases with the variability of the historical log-likelihood.
Since the historical likelihood factorizes over observations, we can write
\begin{equation}
    \log L_0(D_0\mid\theta)
    =
    \sum_{i=1}^{N_0}\log L_i(d_{0i}\mid\theta).
\end{equation}
The width of a global envelope for $Z_{\eta,\eta'}=(\eta'-\eta)\log L_0(D_0\mid\theta)$ typically grows with $N_0$.
This motivates a divide-and-conquer construction.

Let $D_{0,1},\ldots,D_{0,G}$ be a partition of the historical data.
Let $\mathcal I_g$ denote the set of observation indices in block $g$.
Define the block likelihood by
\begin{equation}
    L_{0,g}(D_{0,g}\mid\theta)
    =
    \prod_{i\in\mathcal I_g} L_i(d_{0i}\mid\theta).
\end{equation}
Then the full historical likelihood satisfies
\begin{equation}
    L_0(D_0\mid\theta)
    =
    \prod_{g=1}^{G} L_{0,g}(D_{0,g}\mid\theta).
\end{equation}
Equivalently, the full historical log-likelihood satisfies
\begin{equation}
    \log L_0(D_0\mid\theta)
    =
    \sum_{g=1}^{G}\log L_{0,g}(D_{0,g}\mid\theta).
\end{equation}

For each block $g$, define the blockwise random variable
\begin{equation}
    Z_{g,\eta,\eta'}(U,\theta)
    =
    (\eta'-\eta)\log L_{0,g}(D_{0,g}\mid\theta).
\end{equation}
Assume that deterministic bounds $a_{g,\eta,\eta'}$ and $b_{g,\eta,\eta'}$ are available such that
\begin{equation}
    a_{g,\eta,\eta'}
    \le
    Z_{g,\eta,\eta'}(U,\theta)
    \le
    b_{g,\eta,\eta'}
    \qquad
    \mathsf P_{\eta,\eta'}\text{-almost surely}.
\end{equation}
Set
\begin{equation}
    c_{g,\eta,\eta'}
    =
    b_{g,\eta,\eta'},
    \qquad
    \lambda_{g,\eta,\eta'}
    =
    b_{g,\eta,\eta'}-a_{g,\eta,\eta'}.
\end{equation}
A blockwise Poisson estimator can then be constructed as
\begin{equation}
    \widehat I_g(\eta,\eta')
    =
    \exp\{\lambda_{g,\eta,\eta'}-c_{g,\eta,\eta'}\}
    \prod_{j=1}^{K_g}
    \frac{
        c_{g,\eta,\eta'} -
        Z_{g,\eta,\eta'}(U_{g,j},\theta_{g,j})
    }{
        \lambda_{g,\eta,\eta'}
    },
\end{equation}
where $K_g\sim\mathrm{Poisson}(\lambda_{g,\eta,\eta'})$.
The variables $(U_{g,j},\theta_{g,j})$ are sampled independently from the same proposal-specific measure $\mathsf P_{\eta,\eta'}$.
The blockwise estimator satisfies
\begin{equation}
    \bE[\widehat I_g(\eta,\eta')]
    =
    \exp\{-\Delta_g(\eta,\eta')\},
\end{equation}
where
\begin{equation}
    \Delta_g(\eta,\eta')
    =
    \bE_{\mathsf P_{\eta,\eta'}}
    \left[
        Z_{g,\eta,\eta'}(U,\theta)
    \right].
\end{equation}

\begin{proof}
If the block estimators are independent, the divide-and-conquer product estimator is
\begin{equation}
    \widehat I_{\mathrm{dc}}(\eta,\eta')
    =
    \prod_{g=1}^{G}\widehat I_g(\eta,\eta').
\end{equation}
By independence, its expectation is
\begin{align}
    \bE[\widehat I_{\mathrm{dc}}(\eta,\eta')]
    &=
    \prod_{g=1}^{G}
    \bE[\widehat I_g(\eta,\eta')] \\
    &=
    \prod_{g=1}^{G}
    \exp\{-\Delta_g(\eta,\eta')\} \\
    &=
    \exp\left\{
        -\sum_{g=1}^{G}\Delta_g(\eta,\eta')
    \right\}.
\end{align}
Since
\begin{equation}
    \sum_{g=1}^{G}\Delta_g(\eta,\eta')
    =
    \Delta(\eta,\eta'),
\end{equation}
we obtain
\begin{equation}
    \bE[\widehat I_{\mathrm{dc}}(\eta,\eta')]
    =
    \exp\{-\Delta(\eta,\eta')\}.
\end{equation}
Thus, the divide-and-conquer product estimator is unbiased for $c_0(\eta)/c_0(\eta')$.
\end{proof}

The main advantage of this construction is variance and cost control.
Instead of using one global envelope with width $\lambda_{\eta,\eta'}$, the method uses several smaller envelopes with widths $\lambda_{g,\eta,\eta'}$.
The expected number of Poisson factors is
\begin{equation}
    \sum_{g=1}^{G}\lambda_{g,\eta,\eta'}.
\end{equation}
When the blockwise envelopes are much tighter than the full-data envelope, the product construction can be substantially more stable than the global estimator.

This product estimator should not be confused with a genuine recursive divide-and-conquer Bernoulli factory.
The estimator $\widehat I_{\mathrm{dc}}(\eta,\eta')$ is an unbiased estimator of a normalising-constant ratio.
However, using it as a noisy plug-in inside the Barker probability does not generally yield an exact Barker transition.
A true DCBF Barker transition recursively constructs the accept--reject coin itself.
At the leaves, one simulates smaller Bernoulli coins associated with subsets of likelihood factors.
At internal nodes, these coins are merged recursively until a single coin with the Barker acceptance probability is obtained at the root.
Therefore, the recursive DCBF is a coin construction, not merely a product of random ratio estimates.

\section{Pseudo-marginal via importance sampling}
\label{sec:pseudo-marginal-is}

An alternative to constructing an unbiased estimator of the ratio $c_0(\eta)/c_0(\eta')$ is to directly estimate the normalising constant $c_0(\eta)$ itself and form the ratio from two independent estimates.
% While this does not resolve the Jensen bias when the log-ratio is exponentiated, it provides a practical family of estimators that are straightforward to implement.

The simplest such approach is the \emph{prior Monte Carlo estimator}, which exploits the representation
\begin{align*}
    c_0(\eta) &= \bE_{\theta\sim\pi_0}\qty[L_0(\theta\mid D_0)^\eta].
\end{align*}
Samples drawn directly from the prior $\pi_0$ are cheap to obtain and require no gradient information or auxiliary computations.
The main limitation is that the prior and the power posterior $\pi_\eta(\theta) \propto L_0(\theta \mid D_0)^\eta \pi_0(\theta)$ can be poorly aligned, particularly when the historical data are informative, leading to high variance in the importance weights and an inefficient estimator.

A more principled alternative is \emph{importance sampling} centered at the maximum likelihood estimate. Writing the normalising constant as
\begin{equation*}
    c_0(\eta) = \int_{\theta\in\Theta} L_0(\theta\mid D_0)^\eta \pi_0(\theta)\dd{\theta} = \bE_{\theta\sim\pi_\text{MLE}}\qty[L_0(\theta\mid D_0)^\eta\dfrac{\pi_0(\theta)}{\pi_\text{MLE}(\theta)}],
\end{equation*}
one draws samples from a proposal $\pi_\text{MLE}$ centered at the MLE and reweights by the ratio $\pi_0(\theta)/\pi_\text{MLE}(\theta)$.
The MLE-centered proposal concentrates mass where the likelihood is large, reducing variance relative to the prior estimator when the historical data are informative. The cost is the need to compute or approximate the MLE, which may itself require numerical optimization.

